\section{Training, and where the framework belongs}
\label{sec:training}

\subsection{The framework boundary}
\label{sec:framework-boundary}

The single most useful structural decision in a Rust implementation is where the
deep-learning framework is allowed to appear. Here it appears in exactly one
place: the crate that owns the network. The rules engine has no dependency on it,
the search has no dependency on it, and the parallelism layer has no dependency
on it. Data crosses the boundary as plain arrays and scalars.

The concrete flow is worth writing out, because three design decisions arrive at
the same diagram:

\begin{lstlisting}[style=plaintext,caption={The only framework-aware seam in the
system. Everything above the dotted line compiles and runs with no accelerator
present.},label={lst:seam}]
board  (rules engine: bitboards, absolute colours)
  |  encode()                          -- no framework types
  v
planes [C, 17, 17] u8   (player-relative view, border ring set)
  |  convert                           -- THE ONLY framework-aware step
  v
Tensor<B, 4> [batch, C, 17, 17]
  |  forward
  v
policy logits [batch, 225] , value [batch, 1]
\end{lstlisting}

Three properties follow, and each is a practical benefit rather than an aesthetic
one.

\emph{The search is testable without a device.} Because the evaluator trait
carries plain data (Section~\ref{sec:evaluator}), the entire search can be built
and verified against a deterministic mock evaluator, on a machine with no
accelerator, before any network code exists. This is the difference between
debugging a search and debugging a search and a network simultaneously.

\emph{Compile times stay usable.} A framework and its backend machinery are heavy
dependencies. Confining them to one crate keeps incremental builds of the
frequently edited crates fast, which matters over a project measured in months.

\emph{The relative view is derived, never stored.} The engine stores absolute
colours, because the Swap2 opening does not alternate players
(Section~\ref{sec:solved}); the encoder produces the player-relative planes,
which halves the state space the network must learn. Mixing the two is a
well-known source of subtle bugs, and separating them physically -- different
crate, different types -- makes mixing them a compile error rather than a
misread.

\paragraph{Planes and the border ring.} A detail of the encoding carries more
learning consequence than its size suggests. The board is $15\times15$, but the
encoded planes are $17\times17$, with an extra ring of cells whose contents are
set to the \emph{opponent} for every position. The reason is convolutional
symmetry: with a border, a same-padding convolution at an edge cell sees a wall
of opponent stones, which is how a wall really constrains a line, instead of
seeing zeros and having to learn what an edge means. The border is a property of
the encoding layer and not of the engine's storage, so the engine keeps
$15\times15$ semantics throughout. The policy head still emits 225 logits, one
per real cell, so border cells can never become move targets.

One invariant makes this safe, and it is testable: the border ring is invariant
under the board's dihedral group, and the symmetry transforms are defined on the
$15\times15$ board, so encoding commutes with every transform. Applying a
transform and then encoding must equal encoding and then applying the transformed
transform. That property test catches an entire class of augmentation bugs,
including the ones that only affect one of the eight transforms.

\subsection{Backend genericity, and why it is not ceremony}
\label{sec:backends}

Training code should be generic over an autodiff backend; inference code should
be generic over a plain backend. In Burn's type system that means a model written
as \code{Module<B: Backend>}, instantiated during training with
\code{Autodiff<B>} to obtain gradients and during inference with \code{B} alone.
The binary fixes the concrete choice -- on a CPU-only path, the pure-Rust
\code{Flex} backend; with an accelerator feature enabled, the GPU backend.

Generic backends are usually presented as portability. The real benefit here is
\emph{testing}. Two gates depend on it directly.

The \emph{overfit test} is cheap on a CPU backend and proves the loss, the
encoding, and the training loop are wired together correctly: take a handful of
positions, train on them alone, and require the policy accuracy to approach
certainty. A loop that cannot overfit a few dozen samples has a defect that no
amount of training data will fix.

The \emph{backend parity test} is the executable form of a genuine risk. The same
weights and the same batch, evaluated on the CPU backend and on the accelerator
backend, must agree to within a small tolerance on the forward pass -- and, more
importantly, one full training step must produce parameter deltas within a
relative tolerance. Forward agreement alone does not establish that gradients or
optimizer updates are correct on the accelerator. Until the parity test passes,
training belongs on the CPU backend; once it passes, accelerator training is a
measured claim rather than an assumption. Deterministic-operation concerns on
accelerators, including non-deterministic autotuning, are the reason the
tolerance is a tolerance and not an equality.

\subsection{The training loop}
\label{sec:loop}

The AlphaZero training step needs a manual loop rather than a supervised-training
helper, for two reasons that are properties of the data and not of the
framework. Samples do not exist until they are generated: a mini-batch is
assembled by drawing games from the replay window, replaying each to a random
ply, encoding, and applying augmentation, all during training. And the policy
target is a full distribution over a masked action set, not a class index. Both
are ordinary in AlphaZero and both are outside the shape a supervised trainer
expects.

The loop is therefore: draw a game, replay it to a ply, encode, augment, and
repeat until the batch is full; forward; compute the loss; backward; step. The
only subtlety worth naming is that the augmentation transform must be applied to
the input planes \emph{and} to the policy target consistently, and that a move
index must be remapped through the same transform. Applying the transform to the
board but not to the target distribution is a silent corruption: the loss still
decreases, and the network learns a consistently permuted policy.

\subsection{The loss in practice}
\label{sec:loss-practice}

Equation \eqref{eq:alphazero-loss} becomes three pieces of ordinary code.

The value term is a mean squared error between the value head's output and the
scalar $z$, which takes values in $\{-1,0,+1\}$.

The policy term is a cross-entropy against a \emph{soft} target, which is the one
place a standard loss function cannot be used directly, because most
cross-entropy implementations take a class index. The expression is

\begin{equation}
\mathcal{L}_{\text{pol}}(\pi, p) \;=\; - \sum_{a \in \mathcal{A}_{\text{legal}}}
\pi_a \log p_a ,
\end{equation}

where $p$ is the softmax of the network's logits over the legal actions. Compute
it as \code{-(pi * log\_softmax(logits)).sum()} for numerical stability, and
restrict the sum to legal actions, which is the same masking discipline as
Section~\ref{sec:expand}: the target $\pi$ is already zero on illegal cells, but
an explicit mask keeps a bug in the target construction from silently leaking
mass.

The regularization term is normally not written in the loss at all but expressed
through the optimizer, which brings the next decision.

\subsection{Optimizer and schedule}
\label{sec:optimizer}

The primary sources train with stochastic gradient descent and momentum $0.9$,
with an $L_2$ penalty of $10^{-4}$ inside the loss and a step-decay learning-rate
schedule \citep{silver2017alphagozero,silver2018alphazero}.
Neither paper explains the choice, which is itself informative: in 2017 the
recipe for a deep residual network was SGD with momentum and step decay, and the
architecture was chosen to match the state of the art in image recognition
\citep{he2016resnet}. The optimizer came with the architecture, and every major
replication kept it \citep{tian2019elf,wu2019katago}.

For a single-device run at a much smaller batch size, AdamW
\citep{kingma2014adam,loshchilov2017adamw} is a defensible substitute, and the
reason is worth stating precisely so that the substitution is not mistaken for a
free improvement. With SGD, an $L_2$ penalty in the loss and decoupled weight
decay are nearly the same mechanism, differing by a learning-rate factor. With
Adam they are not the same at all, because the penalty's gradient passes through
the adaptive preconditioner; that difference is the AdamW result. So the
mechanism the papers rely on -- a weight-decay rate of $10^{-4}$ acting as a
smoothness control on a search-critical function
(Section~\ref{sec:loss}) -- can be preserved under AdamW by using decoupled
decay. The reasons to prefer adaptive methods at small scale are the standard
ones: faster early convergence and far less sensitivity to the learning rate.
The reason the papers did not need them is the standard counter-argument:
adaptive methods have been reported to generalize worse on image-classification-
style tasks, and training on a distribution that never stops moving is a
generalization problem throughout the run \citep{wilson2017adaptive}.

Both choices should be configuration, not a rewrite, so that a fidelity
experiment is a one-line change.

The learning-rate schedule needs the same treatment. AlphaGo Zero stepped
$10^{-2} \to 10^{-3} \to 10^{-4}$; AlphaZero started at $0.2$ and dropped three
times, in step with a doubled batch size. Those schedules are tuned for SGD at
batches of two to four thousand and do not transfer to a batch of a thousand with
an adaptive optimizer. A linear warmup followed by cosine decay is the more
useful starting point, and the specific numbers are an experiment, not a
citation.

\subsection{Records, checkpoints, and one trap}
\label{sec:records}

Checkpointing a model in a reduced-precision format is a standard space
optimization and a bad idea in this system. If candidate and baseline are
compared by playing games against each other, and the candidate's weights have
been rounded to half precision while the baseline's have not, the comparison
measures quantization noise as well as training. Precision loss of that kind is
exactly the size of the differences the arena is trying to resolve. The rule is
therefore: store models at full precision.

Checkpoint often, and keep the checkpoints. A run measured in days should never
lose more than one iteration to a fault, and disk is cheaper than a week of
compute.

\subsection{Verification gates}
\label{sec:gates}

The order of the following gates matters, because each one makes the next one
meaningful.

\begin{enumerate}
  \item \textbf{Rules engine, differentially.} A fast bitboard implementation
        against a naive array implementation, on ten thousand random games,
        compared ply by ply. Any disagreement is a defect in one of them.
  \item \textbf{Search, against known facts.} Tactical puzzles with verifiable
        solutions, driven by the deterministic mock evaluator: a win in one must
        be found at a small simulation budget, forced blocks must be played,
        short forced sequences must be found. This is the test that catches the
        sign convention of Section~\ref{sec:backup}.
  \item \textbf{Prover soundness.} Every emitted line is re-verified by replaying
        it while enumerating all defender alternatives. On shallow positions,
        brute-force adjudication must agree with every proof claim. The prover
        may miss wins; it must never certify one that is not there.
  \item \textbf{Network shapes and overfitting.} Forward shapes on both backends,
        and the overfit test of Section~\ref{sec:backends}.
  \item \textbf{Backend parity.} One full training step, both backends,
        parameter deltas within tolerance.
\end{enumerate}

Nothing trains at length until gate 5 passes. The reason is that the failure it
detects -- a correct-looking forward pass with wrong gradients -- produces a run
that looks normal, costs days, and produces a weak agent. That failure is
indistinguishable from a bad hyperparameter at the point where it is observed,
which is why it must be excluded early and mechanically.
